\documentclass{report}
\usepackage{amsfonts, amsmath}
\newcommand\R{{\mathbb R}}
\newcommand\Q{{\mathbb Q}}
\newcommand\N{{\mathbb N}}
\newcommand\Z{{\mathbb Z}}
\pagestyle{empty}
\begin{document}
\large
\centerline{\Huge Analisi Matematica II}
\renewcommand{\thefootnote}{ } 
\footnote{\sl Avete 3 ore di tempo. Ogni esercizio
vale sei punti. La sufficienza si ottiene con un punteggio $\geq$
18. Solo le risposte chiaramente giustificate saranno prese in
considerazione.}
\renewcommand{\thefootnote}{\arabic{footnote}} 
\vskip.1cm
\centerline{\Large Esonero 03-05-2002}
\vskip1cm
\begin{enumerate}
\item Si consideri la funzione $g:\R\to\R$ definita da $g(x):=2x$.
Per ogni funzione $f:\R\to\R$ periodica di periodo $2\pi$ (cio\`e
$f(x+2\pi n)=f(x)$ per ogni $x\in\R$ e $n\in\Z$) si dimostri che $f\circ g$ \`e
anche essa periodica di periodo $2\pi$. Si calcoli la serie di Fourier
di $f\circ g$ in termini di quella di $f$. Dimostrare che esiste una
sola funzione differenziabile tale che $f\circ g=f$ e $f(0)=2$.

\item Si risolva il seguente problema di Cauchy
\begin{align}\label{es1}
       &y'(x)=xy(x)\\
       &y(0)=1
\end{align}
\item Dimostrare che la soluzione di \eqref{es1} \`e differenziabile un
numero infinito di volte e le sue derivate soddisfano le equazioni
\[
y^{(n+1)}(x)=xy^{(n)}(x)+ny^{(n-1)}(x)
\]
per ogni $n\in\N\cup\{0\}$. Usare tale informazione per scrivere la
serie di Taylor di $y(x)$ nel punto $x=0$.

\item Si calcoli il seguente integrale
\[
\int_0^1\frac{x+1}{4-x^2}dx.
\]
\item Si studi la convergenza della serie di funzioni
\[
\sum_{n=1}^\infty\left[ \sqrt{1+\frac{x}{n^2}}-1\right]
\]
per $x\in[0,1]$.
\item
Si calcoli la derivata della funzione
\[
F(x)=\int_0^{x^3}e^{-y^2} y^3 dy.
\]
\end{enumerate}
\end{document}
